\section{理解 Transformer}

\subsection{完整架构}

Transformer 由 encoder 和 decoder 两部分组成。encoder 反复堆叠 self-attention 与 feedforward，decoder 则在 self-attention 之外，再增加一层 cross-attention 去读取 encoder 输出。

\begin{figure}[H]
\centering
\includegraphics[width=0.95\textwidth]{slides-images/slide-026.jpg}
\caption{Transformer encoder-decoder 的完整图：输入嵌入、位置编码、encoder 堆叠、decoder 堆叠、线性层和 softmax\protect\footnotemark}
\end{figure}
\footnotetext{Slides 第27页。}

\subsection{Self-Attention 的直觉}

Attention 可以被理解成一个 ``fuzzy hashtable''：不是精确查找一个 key，而是根据 query 与所有 key 的匹配程度，返回一个加权和。

\begin{figure}[H]
\centering
\includegraphics[width=0.95\textwidth]{slides-images/slide-028.jpg}
\caption{Attention 的直觉：query 与多个 key 做匹配，再对 value 做加权求和\protect\footnotemark}
\end{figure}
\footnotetext{Slides 第29页。}

\begin{importantbox}{self-attention 的计算步骤}
\begin{enumerate}
    \item 对每个词计算 query、key、value
    \item 用 query 和所有 key 计算 attention score
    \item 对 score 做 softmax 归一化
    \item 对 value 做加权求和，得到该词的新表示
\end{enumerate}
\end{importantbox}

\subsection{从公式到向量化}

对第 $i$ 个位置，令
\[
q_i = W_Q x_i,\quad k_i = W_K x_i,\quad v_i = W_V x_i.
\]
则 attention 权重为
\[
\alpha_{ij} = \mathrm{softmax}_j\!\left(\frac{q_i^\top k_j}{\sqrt{d_k}}\right),
\]
输出为
\[
o_i = \sum_j \alpha_{ij} v_j.
\]
把整个序列堆起来后，向量化实现更简单：
\[
Q = XW_Q,\quad K = XW_K,\quad V = XW_V.
\]

\begin{figure}[H]
\centering
\includegraphics[width=0.95\textwidth]{slides-images/slide-029.jpg}
\caption{Transformer encoder 中 self-attention 的逐步计算方法\protect\footnotemark}
\end{figure}
\footnotetext{Slides 第30页。}

\begin{figure}[H]
\centering
\includegraphics[width=0.95\textwidth]{slides-images/slide-030.jpg}
\caption{向量化 self-attention：把整个序列作为矩阵输入一次性计算\protect\footnotemark}
\end{figure}
\footnotetext{Slides 第31页。}

\subsection{为什么 attention 还不够}

如果只有 self-attention，模型本质上只是把 value 向量重新加权平均。这样虽然强大，但表达力仍然不够，所以 Transformer 在每个 attention block 后面还加了 feedforward 层。

\begin{figure}[H]
\centering
\includegraphics[width=0.95\textwidth]{slides-images/slide-032.jpg}
\caption{仅有 self-attention 的问题：缺少逐元素非线性，容易退化为重新混合 value 向量\protect\footnotemark}
\end{figure}
\footnotetext{Slides 第33页。}

\subsection{Residual、LayerNorm 和 Scaled Dot-Product}

为了让深层 Transformer 可训练，原论文引入了三类训练技巧。

\begin{figure}[H]
\centering
\includegraphics[width=0.95\textwidth]{slides-images/slide-033.jpg}
\caption{训练 Transformer 的三类关键技巧：residual connection、layer norm、scaled dot-product attention\protect\footnotemark}
\end{figure}
\footnotetext{Slides 第34页。}

\begin{knowledgebox}{Residual Connection}
残差连接把输入直接加到子层输出上：
\[
x' = x + \mathrm{Sublayer}(x).
\]
它的作用是保留原始信息、帮助梯度传播，并让深层网络更容易学习到接近恒等映射的行为。
\end{knowledgebox}

\begin{knowledgebox}{Layer Normalization}
LayerNorm 在每个 token 的特征维度上做归一化，减少 layer 输入分布的漂移：
\[
\mathrm{LN}(x) = \gamma \odot \frac{x - \mu}{\sqrt{\sigma^2 + \epsilon}} + \beta.
\]
它比 batch norm 更适合 NLP，因为句子长度变化大，batch 的统计量也不稳定。
\end{knowledgebox}

\begin{knowledgebox}{Scaled Dot-Product Attention}
当 $d_k$ 很大时，$q^\top k$ 的方差会变大，softmax 容易饱和。于是要除以 $\sqrt{d_k}$，把数值尺度控制在更稳定的范围内。
\end{knowledgebox}

\begin{figure}[H]
\centering
\includegraphics[width=0.95\textwidth]{slides-images/slide-034.jpg}
\caption{Residual connections：把原始输入直接传到下一层，有助于训练深层网络\protect\footnotemark}
\end{figure}
\footnotetext{Slides 第35页。}

\begin{figure}[H]
\centering
\includegraphics[width=0.95\textwidth]{slides-images/slide-035.jpg}
\caption{LayerNorm：在每层内部归一化，稳定训练分布\protect\footnotemark}
\end{figure}
\footnotetext{Slides 第36页。}

\begin{figure}[H]
\centering
\includegraphics[width=0.95\textwidth]{slides-images/slide-037.jpg}
\caption{缩放点积注意力：除以 $\sqrt{d_k}$ 后，dot product 的尺度更稳定\protect\footnotemark}
\end{figure}
\footnotetext{Slides 第38页。}

\subsection{位置编码：解决顺序信息缺失}

self-attention 本身不关心词序。若直接把词向量送进去，``Man eats small dinosaur'' 和 ``Dinosaur eats small man'' 的表示几乎没有顺序差别，这显然不对。因此需要把位置编码注入模型。

\begin{figure}[H]
\centering
\includegraphics[width=0.95\textwidth]{slides-images/slide-040.jpg}
\caption{通过 positional encoding 把顺序信息注入输入表示\protect\footnotemark}
\end{figure}
\footnotetext{Slides 第41页。}

一种经典做法是把位置向量 $p_i$ 加到 token 表示上：
\[
x_i' = x_i + p_i.
\]
原始 Transformer 使用正弦/余弦位置编码：
\[
p_{i,2k} = \sin\!\left(i / 10000^{2k/d}\right), \quad
p_{i,2k+1} = \cos\!\left(i / 10000^{2k/d}\right).
\]
这种设计的优点是可外推到更长的序列，而且不同频率的正弦信号能编码相对位置信息。

\begin{figure}[H]
\centering
\includegraphics[width=0.95\textwidth]{slides-images/slide-041.jpg}
\caption{正弦位置表示：用不同周期的 sin/cos 组合表示绝对位置\protect\footnotemark}
\end{figure}
\footnotetext{Slides 第42页。}

\begin{figure}[H]
\centering
\includegraphics[width=0.95\textwidth]{slides-images/slide-042.jpg}
\caption{Relative position encoding：有时词与词的相对关系比绝对位置更重要\protect\footnotemark}
\end{figure}
\footnotetext{Slides 第43页。}

\subsection{多头注意力}

单个 attention head 往往只能学到一种局部模式。多头注意力让模型并行地学习多个投影空间里的匹配关系。

\begin{figure}[H]
\centering
\includegraphics[width=0.95\textwidth]{slides-images/slide-044.jpg}
\caption{多头 self-attention：k 个 head 并行工作，最后把结果拼接起来\protect\footnotemark}
\end{figure}
\footnotetext{Slides 第45页。}

若有 $h$ 个头，第 $\ell$ 个头使用自己的参数：
\[
Q^{(\ell)}, K^{(\ell)}, V^{(\ell)} \in \mathbb{R}^{d \times d/h}.
\]
每个头独立计算 attention，再把输出拼接：
\[
\mathrm{MHAtt}(X) = Y [O^{(1)}; \dots; O^{(h)}].
\]
直观上，不同的 head 可以分别关注词法邻近、句法依赖、长程共指等不同模式。

\subsection{Decoder 的两类 attention}

Transformer decoder 需要先做 masked self-attention，再做 encoder-decoder attention。前者防止“偷看未来”，后者从 encoder 的 memory 中读取源句信息。

\begin{figure}[H]
\centering
\includegraphics[width=0.95\textwidth]{slides-images/slide-046.jpg}
\caption{完成 encoder 后进入 decoder 部分\protect\footnotemark}
\end{figure}
\footnotetext{Slides 第47页。}

\begin{figure}[H]
\centering
\includegraphics[width=0.95\textwidth]{slides-images/slide-047.jpg}
\caption{Masked self-attention 的问题：如果不遮挡未来 token，decoder 会直接作弊\protect\footnotemark}
\end{figure}
\footnotetext{Slides 第48页。}

\begin{figure}[H]
\centering
\includegraphics[width=0.95\textwidth]{slides-images/slide-048.jpg}
\caption{解决方法：在 attention score 上对未来位置加 mask，让它们变成 $-\infty$\protect\footnotemark}
\end{figure}
\footnotetext{Slides 第49页。}

\[
e_{ij} =
\begin{cases}
q_i^\top k_j, & j < i, \\
-\infty, & j \ge i.
\end{cases}
\]
这保证了训练时可以并行计算所有位置，但每个位置仍然只能依赖左侧上下文。

\begin{figure}[H]
\centering
\includegraphics[width=0.95\textwidth]{slides-images/slide-050.jpg}
\caption{Decoder 的 masked multi-head self-attention 模块\protect\footnotemark}
\end{figure}
\footnotetext{Slides 第51页。}

\begin{figure}[H]
\centering
\includegraphics[width=0.95\textwidth]{slides-images/slide-051.jpg}
\caption{Encoder-decoder attention：query 来自 decoder，keys 和 values 来自 encoder\protect\footnotemark}
\end{figure}
\footnotetext{Slides 第52页。}

\begin{knowledgebox}{Self-attention 与 cross-attention 的区别}
\begin{itemize}
    \item \textbf{Self-attention}：$Q,K,V$ 都来自同一序列
    \item \textbf{Cross-attention}：$Q$ 来自 decoder，$K,V$ 来自 encoder
\end{itemize}
前者负责建模句内依赖，后者负责在生成时读取输入句子的记忆。
\end{knowledgebox}

\subsection{Decoder 的收尾}

decoder 最后还会经过 feedforward、线性层和 softmax，变成下一个 token 的概率分布。

\begin{figure}[H]
\centering
\includegraphics[width=0.95\textwidth]{slides-images/slide-052.jpg}
\caption{Decoder 再次叠加 attention、feedforward、residual 和 layer norm\protect\footnotemark}
\end{figure}
\footnotetext{Slides 第53页。}

\begin{figure}[H]
\centering
\includegraphics[width=0.95\textwidth]{slides-images/slide-054.jpg}
\caption{输出侧的最终线性映射：把 hidden state 投影到词表维度\protect\footnotemark}
\end{figure}
\footnotetext{Slides 第55页。}

\begin{figure}[H]
\centering
\includegraphics[width=0.95\textwidth]{slides-images/slide-055.jpg}
\caption{最后接 softmax，得到下一词的概率分布\protect\footnotemark}
\end{figure}
\footnotetext{Slides 第56页。}

\subsection{本章小结}

Transformer 的主体可以概括为：\textbf{位置编码 + 多头 self-attention + 残差连接 + layer norm + feedforward}。encoder 和 decoder 的差异主要在于 decoder 多了 mask 和 cross-attention。

%% ========== 缺点与变体 ========== 

